% GENERATED FILE. Edit canonical lessons, then run npm run generate:books.
% canonical-source-hash: fnv1a64:a63103d8c1eef295
% canonical-lessons: TE-C406-aviniti, TE-C406-lancam, TE-C406-paradarsakata, TE-C406-aropana, TE-C406-aropincu

\chapter{Corruption and Accountability}
\label{ch:te-corruption-and-accountability}
\hlchaptermodality{406}

\begin{chapteropening}
\textbf{By the end of this chapter:} \emph{I can talk about corruption, bribes and allegations.}

Complete the last lesson of chapter 406: I can talk about corruption, bribes and allegations.
\end{chapteropening}

\section[\texorpdfstring{\te{అవినీతి} (avinīti)}{అవినీతి (avinīti)}]{\te{అవినీతి} (avinīti) — corruption}

\label{lesson:TE-C406-aviniti}



\begin{quote}
Before the new one: say the Telugu for a budget, then the Telugu for a loan.
\end{quote}

\subsection*{You'll want to know: \te{అవినీతి}}
\textbf{\te{అవినీతి}} — \emph{avinīti} — \textquotedblleft{}corruption\textquotedblright{}.

\begin{quote}
\te{అవినీతి} \te{వల్ల} \te{దేశానికి} \te{చాలా} \te{నష్టం} \te{జరుగుతుంది}.
\end{quote}

\begin{quote}
\emph{avinīti valla dēśāniki cālā naṣṭaṁ jarugutundi.}
\end{quote}

\begin{quote}
Corruption does a lot of harm to the country.
\end{quote}

\te{అ}-, not, + \te{వినీతి}, good conduct.

\begin{cousinweb}
\textbf{\te{అవినీతి}} (\emph{avinīti}) is built from Sanskrit \emph{a-}, \textquotedblleft{}not\textquotedblright{}, and \emph{vinīti}, \textquotedblleft{}training, good behaviour\textquotedblright{} (Monier-Williams).
\end{cousinweb}

\begin{grammarlens}[title={What you've built}]
One more word for giving reasons and opinions.
\end{grammarlens}

\subsection*{Your turn}
\begin{itemize}
\raggedright
  \item \emph{Say it:} \emph{avinīti}
  \item \emph{Say it:} \emph{avinīti}, once more
  \item \emph{Say it:} \emph{avinīti valla dēśāniki cālā naṣṭaṁ jarugutundi.} — Corruption does a lot of harm to the country.
  \item \emph{Recall:} say the Telugu for a budget, then the Telugu for a loan, then say \emph{avinīti} again
\end{itemize}

\begin{tcolorbox}[breakable,colback=teal!4,colframe=teal!35!black,title={Before you move on}]
What does \te{అవినీతి} mean? (\textquotedblleft{}Corruption\textquotedblright{}.) Say it once more.
\end{tcolorbox}
\section[\texorpdfstring{\te{లంచం} (lañcaṁ)}{లంచం (lañcaṁ)}]{\te{లంచం} (lañcaṁ) — a bribe}

\label{lesson:TE-C406-lancam}



\begin{quote}
Before the new one: say the Telugu for a loan, then the Telugu for corruption.
\end{quote}

\subsection*{You'll want to know: \te{లంచం}}
\textbf{\te{లంచం}} — \emph{lañcaṁ} — \textquotedblleft{}a bribe\textquotedblright{}.

\begin{quote}
\te{ఆ} \te{అధికారి} \te{లంచం} \te{అడిగాడు}.
\end{quote}

\begin{quote}
\emph{ā adhikāri lañcaṁ aḍigāḍu.}
\end{quote}

\begin{quote}
That officer asked for a bribe.
\end{quote}

Plural: \te{లంచాలు}, \emph{lañcālu}.

\te{లంచం} \te{ఇచ్చు}, to give a bribe; \te{లంచం} \te{తీసుకో}, to take one.

\begin{cousinweb}
\textbf{\te{లంచం}} (\emph{lañcaṁ}) is borrowed from Sanskrit \emph{lañcā}, \textquotedblleft{}a present, a bribe\textquotedblright{} (CDIAL 10914).
\end{cousinweb}

\begin{grammarlens}[title={What you've built}]
One more word for giving reasons and opinions.
\end{grammarlens}

\subsection*{Your turn}
\begin{itemize}
\raggedright
  \item \emph{Say it:} \emph{lañcaṁ}
  \item \emph{Say it:} \emph{lañcaṁ}, once more
  \item \emph{Say it:} \emph{ā adhikāri lañcaṁ aḍigāḍu.} — That officer asked for a bribe.
  \item \emph{Recall:} say the Telugu for a loan, then the Telugu for corruption, then say \emph{lañcaṁ} again
\end{itemize}

\begin{tcolorbox}[breakable,colback=teal!4,colframe=teal!35!black,title={Before you move on}]
What does \te{లంచం} mean? (\textquotedblleft{}A bribe\textquotedblright{}.) Say it once more.
\end{tcolorbox}
\section[\texorpdfstring{\te{పారదర్శకత} (pāradarśakata)}{పారదర్శకత (pāradarśakata)}]{\te{పారదర్శకత} (pāradarśakata) — transparency}

\label{lesson:TE-C406-paradarsakata}



\begin{quote}
Before the new one: say the Telugu for corruption, then the Telugu for a bribe.
\end{quote}

\subsection*{You'll want to know: \te{పారదర్శకత}}
\textbf{\te{పారదర్శకత}} — \emph{pāradarśakata} — \textquotedblleft{}transparency\textquotedblright{}.

\begin{quote}
\te{ప్రభుత్వ} \te{పనుల్లో} \te{పారదర్శకత} \te{ఉండాలి}.
\end{quote}

\begin{quote}
\emph{prabhutva panullō pāradarśakata uṇḍāli.}
\end{quote}

\begin{quote}
There should be transparency in the work of government.
\end{quote}

\te{పార}, the far side, + \te{దర్శక}, showing, + -\te{త}.

\begin{cousinweb}
\textbf{\te{పారదర్శకత}} (\emph{pāradarśakata}) is built on Sanskrit \emph{pāradarśaka}, \textquotedblleft{}showing the opposite shore\textquotedblright{} (Monier-Williams), with the abstract ending -\te{త}; the sense \textquotedblleft{}transparency\textquotedblright{} is modern.
\end{cousinweb}

\begin{grammarlens}[title={What you've built}]
One more word for giving reasons and opinions.
\end{grammarlens}

\subsection*{Your turn}
\begin{itemize}
\raggedright
  \item \emph{Say it:} \emph{pāradarśakata}
  \item \emph{Say it:} \emph{pāradarśakata}, once more
  \item \emph{Say it:} \emph{prabhutva panullō pāradarśakata uṇḍāli.} — There should be transparency in the work of government.
  \item \emph{Recall:} say the Telugu for corruption, then the Telugu for a bribe, then say \emph{pāradarśakata} again
\end{itemize}

\begin{tcolorbox}[breakable,colback=teal!4,colframe=teal!35!black,title={Before you move on}]
What does \te{పారదర్శకత} mean? (\textquotedblleft{}Transparency\textquotedblright{}.) Say it once more.
\end{tcolorbox}
\section[\texorpdfstring{\te{ఆరోపణ} (ārōpaṇa)}{ఆరోపణ (ārōpaṇa)}]{\te{ఆరోపణ} (ārōpaṇa) — an allegation; a charge}

\label{lesson:TE-C406-aropana}



\begin{quote}
Before the new one: say the Telugu for a bribe, then the Telugu for transparency.
\end{quote}

\subsection*{You'll want to know: \te{ఆరోపణ}}
\textbf{\te{ఆరోపణ}} — \emph{ārōpaṇa} — \textquotedblleft{}an allegation; a charge\textquotedblright{}.

\begin{quote}
\te{అతని} \te{మీద} \te{అవినీతి} \te{ఆరోపణలు} \te{వచ్చాయి}.
\end{quote}

\begin{quote}
\emph{atani mīda avinīti ārōpaṇalu vaccāyi.}
\end{quote}

\begin{quote}
Allegations of corruption were made against him.
\end{quote}

Plural: \te{ఆరోపణలు}, \emph{ārōpaṇalu}.

The person accused takes \te{మీద}: \te{అతని} \te{మీద} \te{ఆరోపణ}.

\begin{cousinweb}
\textbf{\te{ఆరోపణ}} (\emph{ārōpaṇa}) is borrowed from Sanskrit \emph{āropaṇa}, \textquotedblleft{}placing or fixing on\textquotedblright{} (Monier-Williams), of the family of \emph{āropa}, \textquotedblleft{}charging with, attributing to\textquotedblright{}.
\end{cousinweb}

\begin{grammarlens}[title={What you've built}]
One more word for giving reasons and opinions.
\end{grammarlens}

\subsection*{Your turn}
\begin{itemize}
\raggedright
  \item \emph{Say it:} \emph{ārōpaṇa}
  \item \emph{Say it:} \emph{ārōpaṇa}, once more
  \item \emph{Say it:} \emph{atani mīda avinīti ārōpaṇalu vaccāyi.} — Allegations of corruption were made against him.
  \item \emph{Recall:} say the Telugu for a bribe, then the Telugu for transparency, then say \emph{ārōpaṇa} again
\end{itemize}

\begin{tcolorbox}[breakable,colback=teal!4,colframe=teal!35!black,title={Before you move on}]
What does \te{ఆరోపణ} mean? (\textquotedblleft{}An allegation; a charge\textquotedblright{}.) Say it once more.
\end{tcolorbox}
\section[\texorpdfstring{\te{ఆరోపించు} (ārōpiñcu)}{ఆరోపించు (ārōpiñcu)}]{\te{ఆరోపించు} (ārōpiñcu) — to allege; to accuse}

\label{lesson:TE-C406-aropincu}



\begin{quote}
Before the new one: say the Telugu for transparency, then the Telugu for an allegation.
\end{quote}

\subsection*{You'll want to know: \te{ఆరోపించు}}
\textbf{\te{ఆరోపించు}} — \emph{ārōpiñcu} — \textquotedblleft{}to allege; to accuse\textquotedblright{}.

\begin{quote}
\te{అతను} \te{డబ్బు} \te{తీసుకున్నాడని} \te{ఆమె} \te{ఆరోపించింది}.
\end{quote}

\begin{quote}
\emph{atanu ḍabbu tīsukunnāḍani āme ārōpiñcindi.}
\end{quote}

\begin{quote}
She alleged that he had taken the money.
\end{quote}

Past (I …): \te{ఆరోపించాను}, \emph{ārōpiñcānu}.

What is alleged is a clause ending in \te{అని}.

\begin{cousinweb}
\textbf{\te{ఆరోపించు}} (\emph{ārōpiñcu}) is built on Sanskrit \emph{āropa}, \textquotedblleft{}charging with; assigning or attributing to\textquotedblright{} (Monier-Williams), with the Telugu verb ending -\te{ఇంచు}.
\end{cousinweb}

\begin{grammarlens}[title={What you've built}]
One more word for giving reasons and opinions.
\end{grammarlens}

\subsection*{Your turn}
\begin{itemize}
\raggedright
  \item \emph{Say it:} \emph{ārōpiñcu}
  \item \emph{Say it:} \emph{ārōpiñcu}, once more
  \item \emph{Say it:} \emph{atanu ḍabbu tīsukunnāḍani āme ārōpiñcindi.} — She alleged that he had taken the money.
  \item \emph{Recall:} say the Telugu for transparency, then the Telugu for an allegation, then say \emph{ārōpiñcu} again
\end{itemize}

\begin{tcolorbox}[breakable,colback=teal!4,colframe=teal!35!black,title={Before you move on}]
What does \te{ఆరోపించు} mean? (\textquotedblleft{}To allege; to accuse\textquotedblright{}.) Say it once more.
\end{tcolorbox}
